\scanpage{16}
Transformant l'équation
\[
x^2=y^3+axy+by^2+cx+dy+e
\]
par $s$, \deleted{et} se rappellant que $sx=-x$, $sy=y$, on trouve
\[
x^2=y^3\deleted{\unsure{+}}-axy+by^2-cx+dy+e
\]
et en comparant on trouve $a=c=0$, donc l'équation s'écrit maintenant sous la forme
\[
\boxed{x^2=y^3+ay^2+by+c}
\]
avec comme condition $\delta(a,b,c)\ne0$, qui ici ne fait qu'exprimer que les solutions de $P(y)=y^3+ay^{\unsure{2|1}}+by+c$ sont distinctes en tt pt de $S$.

Géométriquement, \deleted{\unsure{la situ.}} l'équation s'interprète ainsi : Le quotient $C/\mathbb Z_{/2}$ est un schéma de Brauer-Severi \inserted{$P$} sur $S$, muni d'une section marquée (correspondant à la section nulle de $C$ sur $S$), donc \deleted{s'identifie} \deleted{est isomorphe} \inserted{canoniquement} à un espace principal homogène sous un \deleted{fibré \unsure{qui est loc libre}} \deleted{\unsure{faisceau inversible}} fibré vectoriel \deleted{\unsure{homogène}} \inserted{inversible} [il correspond : $L^{\otimes2}$]
% slanted
\margin{\unsure{précisons|préciser|provisoire}}

$C$ est ramifiée \deleted{\unsure{et}} au dessus d'un diviseur relatif \inserted{de $P/S$} étale sur $S$, qui se décompose en la section à l'infini, et en un diviseur \deleted{\unsure{de}} relatif \inserted{de $\unsure{\mathring P|\dot P|P}/S$} étale sur $S$, \deleted{\unsure{et en fait de}} de \deleted{$\unsure{\mathcal E}$} $\underline V(L^{\otimes2})/S$ étale sur $S$, de degré relatif 3.
% archive sheet number: 15
